\begin{abstract}
An incompletely specified action need not admit a completion on its original
carrier, and its unspecified values need not remain independent. We study
finite two-sorted ground presentations consisting of complete native operation
tables, carrier-compatibility equations, and arbitrary supplied action cells.
For this pure-row fragment, we characterize the carrier congruence forced by
initial semantics as the least quotient on which every row pushout preserves
the distinguished carrier. Row protection is equivalent to homomorphic
coherence of the supplied data and an exact kernel-extension condition.
On the resulting quotient, if a row induces $h:S\to Q$ and
$\kappa=\Cg_Q(\ker h)$, its forced domain is precisely the saturation
$S^\kappa$. We identify the full equality relation on the named carrier and
named action cells and give its exact class counts at proof horizons one and
two. Carrier-valued completion has a different criterion: the carrier
embedding into the row pushout must split. Finite examples show that repairs
need not have a least element, need not be upward closed, and can force
strictly more carrier equality than unrestricted models. Further results
describe the effect of composition, group laws, and one-sorted coupling.
The proofs use unrestricted ground-diagram semantics; no identities of a
chosen variety are imposed on unnamed elements.
\end{abstract}

\noindent\textbf{Keywords.} Initial algebra; ground equations; incomplete
action; congruence extension; pushout; kernel saturation; quotient repair.
